% ==============================================================================
% sec:case_studies --- metodika případových studií
% Zdroje: PennyMarket_KS2026, KS_CeskaPosta_2025, KSVS_OSPZV_ASO_2026,
%         KSVS_OSPPP_2026, CMKOS_ZpravaKV2025, CMKOS_PrinosKV2025, PAQ_Svarcsystem,
%         zakon_kv_1991, zakon_zp_2006
% ==============================================================================

Případové studie navazují na~komparativní rovinu zpracovanou v~kap.~\ref{sec:teorie} a~na~empirickou analýzu agregovaných ukazatelů z~kap.~\ref{sec:eu_pokryti_kv}.
Jejich účelem je doložit, jak se obecné rysy českého modelu \aclgen{KV} --- nízká odborová organizovanost, převaha podnikové úrovně vyjednávání a~minimální využití rozšiřování závaznosti \aclgen{KSVS} --- promítají do~konkrétních textů smluv a~do~úrovně skutečně sjednaných pracovních a~mzdových podmínek.\par

Výběr čtyř případů respektuje dvojrozměrnou typologii.
Po~vertikální ose stojí podniková úroveň (\acs{PKS}) proti odvětvové (\acs{KSVS}); po~horizontální ose se~liší sektory podle organizační síly sociálních partnerů a~podle ekonomické pozice odvětví.
\Acf{PKS} \emph{Penny~Market} a~\acs{PKS} \emph{Česká~pošta} reprezentují dva odlišné případy podnikového vyjednávání v~\acloc{ČR} -- soukromý maloobchodní řetězec se~zahraničním vlastníkem a~státní podnik s~koexistencí devíti \aclpgen{OO}.
Na~odvětvové úrovni je analyzována \acs{KSVS} v~zemědělství, jejíž závaznost byla rozšířena Sdělením č.~47/2026~Sb.~na~zaměstnavatele s~převažující činností v~\acgen{NACE}~01.1 až~01.6 a~je tak jednou ze čtyř \acs{KSVS} aktuálně závazných pro další zaměstnavatele.~\cite{sdeleni_47_2026}
Druhým případem je \acs{KSVS} bankovnictví a~pojišťovnictví, jež reprezentuje sektor s~vysokou přidanou hodnotou.\par

Provedená metodika respektuje jednotnou strukturu. Nejprve identifikujeme smluvní strany, vymezení \acgen{NACE} a~platnost k~ověření legitimity dle~\S~22--25 \aclgen{ZP} a~\S~3 \acgen{ZKV}. Následně srovnáváme mzdová minima, příplatky a~další parametry se~zákonnými hranicemi \aclgen{ZP}. Třetím krokem je~vyhodnocení odchylek vůči odvětvovým standardům ze~zpráv \acgen{ČMKOS}.~\cite{CMKOS_ZpravaKV2025}\par

Klíčovým aspektem je~analýza motivací aktérů: kolektivní smlouvy nahlížíme jako reciproční směnu, kde zaměstnavatelé smění určité mzdové či sociální nadstandardy za~flexibilitu organizace práce či garantovaný sociální smír. Tato bilance pomáhá osvětlit výslednou úroveň ochrany v~sektoru.\par

Omezení metodologie případových studií:
\begin{enumerate}
  \item dvě analyzované \acs{PKS} byly vybrány účelově s~ohledem na~kontrastní typologii a~nepředstavují statisticky reprezentativní vzorek,
  \item příloha č.~1 \aclgen{PKS} \emph{Penny~Market} obsahující konkrétní mzdové tarify nebyla v~získaném dokumentu dostupná, takže srovnání mzdových sazeb se~opírá o~ujednanou meziroční změnu a~o~minimální mzdový tarif,
  \item vybrané \acs{PKS} nejsou vztažené k~žádné \acs{KSVS} a~neumožňují tak porovnat vliv odvětvového \acs{KV} na podnikové \ac{KV}.\par
\end{enumerate}
